\documentclass[10pt,letterpaper]{article}
\usepackage[margin=.82in]{geometry}
\usepackage{amsmath,amssymb,amsthm}
\usepackage[hidelinks]{hyperref}
\newtheorem{theorem}{Theorem}
\newcommand{\N}{\mathbb Z_{\geq 0}}
\begin{document}
\begin{center}
{\Large\bfseries The third-largest parts of four-part partitions}\par
\medskip
Pioneer project \quad $\cdot$ \quad 28 November 2026
\end{center}

\begin{abstract}
Let $a_n$ be the sum of the third-largest parts over all partitions of $n$
into four positive parts. We prove the ordinary generating function and
constant-coefficient recurrence proposed by Colin Barker in OEIS A308758.
The proof uses the four gaps of a sorted partition and the linear weight
carried by the two gaps that determine its third-largest part.
\end{abstract}

Write each partition in weakly increasing order $u\leq v\leq w\leq z$, so its contribution to $a_n$ is $v$.
Barker proposed both formulas on June 23, 2019 \cite{oeis}.

\begin{theorem}\label{th:main}
In $\mathbb Z[[x]]$,
\[
\sum_{n\geq0} a_nx^n
 =\frac{x^4(1+x+x^2+x^3+x^4+x^5+x^6)}
 {(1-x)^5(1+x)^3(1+x^2)^2(1+x+x^2)^2}.
\]
For every $n\geq16$,
\begin{align*}
 a_n={}&a_{n-2}+2a_{n-3}+2a_{n-4}-2a_{n-5}-3a_{n-6}-4a_{n-7}\\
       &+4a_{n-9}+3a_{n-10}+2a_{n-11}-2a_{n-12}-2a_{n-13}-a_{n-14}+a_{n-16}.
\end{align*}
\end{theorem}

\begin{proof}
Use the nonnegative gaps $p=u-1$, $q=v-u$, $r=w-v$ and $s=z-w$.
The inverse map has parts
\[
 (1+p,\ 1+p+q,\ 1+p+q+r,\ 1+p+q+r+s).
\]
Thus the total is $4+4p+3q+2r+s$ and the weight is $1+p+q$.
Let $B=(1-x)(1-x^2)(1-x^3)(1-x^4)$. Each quadruple contributes
$(1+p+q)x^{4+4p+3q+2r+s}$. Summing the constant term and the two linear
weights separately gives
\begin{align*}
\sum_{n\geq0} a_nx^n
 &= \sum_{p,q,r,s\in\N}(1+p+q)x^{4+4p+3q+2r+s}\\
 &= \frac{x^4}{B}\left(1+\frac{x^4}{1-x^4}+\frac{x^3}{1-x^3}\right)\\
 &= \frac{x^4(1-x^7)}{(1-x)(1-x^2)(1-x^3)^2(1-x^4)^2}.
\end{align*}
Only finitely many quadruples contribute to each coefficient. The second
equality uses $\sum_{k\geq0}y^k=1/(1-y)$ and
$\sum_{k\geq0}ky^k=y/(1-y)^2$ for $y=x^3$ or $x^4$. Factor $1-x^2$,
$1-x^3$, and $1-x^4$, then cancel the factor $1-x$ from $1-x^7$.
This gives the displayed generating function.

Let $D(x)$ be its denominator. Multiplying the four factors gives
\begin{align*}
D(x)={}&1-x^2-2x^3-2x^4+2x^5+3x^6+4x^7-4x^9-3x^{10}\\
       &-2x^{11}+2x^{12}+2x^{13}+x^{14}-x^{16}.
\end{align*}
The numerator has degree $10$. For $n\geq16$ the coefficient of $x^n$
in $D(x)\sum_{j\geq0}a_jx^j$ therefore vanishes, and every shift in the
resulting equation has a nonnegative index. Solving for $a_n$ gives the recurrence.
\end{proof}

\paragraph{Lean verification.}
The merged Left to Solve statement uses truncated arithmetic sums for the
gaps. The formal proof checks their generating-function rules directly.
For $d>0$, let $A(m)=\sum_{0\leq i\leq m/d}f(m-di)$ and
$W(m)=\sum_{0\leq i\leq m/d}i f(m-di)$, with integer bounds.
Separating $i=0$ gives $A(m+d)=f(m+d)+A(m)$ and $W(m+d)=W(m)+A(m)$.
Together with the initial coefficients, these give
$(1-x^d)\mathcal A(x)=\mathcal F(x)$ and $(1-x^d)\mathcal W(x)=x^d\mathcal A(x)$
for the corresponding ordinary generating functions. The remaining steps are
polynomial identities and a coefficient comparison, checked in Lean 4.33.1.

\begin{thebibliography}{9}
\bibitem{oeis} The OEIS Foundation Inc., \emph{The On-Line Encyclopedia of Integer Sequences},
\href{https://oeis.org/A308758}{entry A308758}.
\end{thebibliography}
\end{document}
